\expected{am-ambassador-v-chicago}
\expected{kam-almaz-v-us}

\item \textbf{Distinguishing cases.} Both of these cases recite essentially the
same rule for formation of a bailment, and the general fact patterns are the
same. How do they reach different results? The easy way out is to point out that
Illinois state law is different from federal law, so there is no need for the
two cases to be consistent with each other. But as clever lawyers, we can do
better.

First, identify the key factual differences between the two cases. Ideally, try
to identify distinctions that seem meaningful and relevant. The fact that
Kam-Almaz's name starts with a ``K'' but Anderson's starts with an ``A'' is a
difference, but there's no good reason why that should change the legal outcome.
What about the fact that one case involved a car and the other a laptop? What
about the nature of the enforcement agencies or the justification for the
seizure?

Next, try to fit the factual distinction into the doctrine. From the rule for
bailment formation, which of the elements are disputed in the above two cases?
Can you make an argument that your factual distinction leads to different
results for the different elements of the rule? Put another way, can you define
the elements of the rule in such a way that makes your distinction relevant?

You might be thinking, ``Why is it \emph{my} job to define the elements of a
rule? Don't courts define these things?'' But remember that, in the common-law
system, rules are interpreted incrementally through the process of cases and
precedents. Legal argumentation in that sense is all about definitions:
Each attorney proposes a preferred interpretation of a legal rule that best
reconciles the case law (and, coincidentally, renders the attorney's client
victorious), and the judge picks the best interpretation to apply.

But what makes one definition better than another, such that it persuades a
court? Of course, one can never actually step into the mind of a judge, and
methodologies for deciding among competing interpretations have been a matter of
debate across the legal academy and profession for centuries. But here is a
simple consideration to think about. Judges, like other humans, want to go to
bed at night thinking they've done good for the world. This is where legal
theory and policy considerations, like those from
\mref{intro-property-rationales}, come in. \emph{Policy and theory are not legal
arguments in themselves, but they justify interpretations of legal rules as
being preferable to other interpretations.}

Put it all together now. Based on the factual differences you find between the
cases, what are plausible ways of defining the elements of bailment formation,
such as delivery and implied contract? And can you argue that those definitions
satisfy some underlying objective of society at large?

\defcase{first-am-v-dc}{
First American Bank v. District of Columbia, 583 A.2d 993 (D.C. 1990)
}

\item \textbf{Bailment injuries.} Once a bailment has been formed, the bailee
owes certain duties to take care of the bailor's property. (Question: does the
bailor owe the bailee any duties?) The standard of care typically depends on the
nature of the relationship between the bailor and bailee. In
\inline{first-am-v-dc}, another case dealing with a towed car, the court deemed
the District to be a bailee and then turned to the duty of care:
\begin{quote}
The question we must resolve is whether the bailment was gratuitous or for hire.
A bailee that takes possession of goods solely for the benefit of the owner is a
gratuitous bailee and liable only for gross negligence, willful acts or fraud.
In contrast, a bailee that receives compensation for its services is
held to a standard of ordinary care [that is, regular negligence]\ldots.
\end{quote}
\sentence{first-am-v-dc at 996}. In this case, the court held that the District
did receive ``compensation'' in the form of (1) ``furthering its interest in
insuring the smooth flow of traffic'' and (2) the fines and fees charged for
return of the vehicle. \sentence{first-am-v-dc at 997}.

\defcase{prime-v-ssi}{
Prime Products, Inc. v. SSI Plastics, Inc., 97 S.W.3d 631 (Tex. App. 2002)
}

Additionally, where the bailee fails to return the property or returns it in
damaged condition, negligence is generally presumed. \sentence{see prime-v-ssi
at 635}.
